% 原书分章PDF第21—24页：21.6模拟全部结果与21.7全部延伸阅读。
\section{模拟结果}

模拟研究基于前述数据生成分布；与前面一样，假定结局 $Y_i$ 是表示在预先指定的随访期内是否维持健身房会员资格的二元指示变量。

我们感兴趣的是考察并估计，对这样的社区施加各种假想干预 $\bm g^*$ 后，平均反事实结局 $E[\bar Y_{\bm g^*}]$ 的平均值。首先，考虑随机干预 $\bm g_1^*$，它以恒定概率0.35将每个个体分配至处理；这不同于观测数据中向25\%的社区成员分配处理的情况。我们还考虑一种情境：上述经济激励受到资源约束，最多只能分配给10\%的社区成员，并估计各种有针对性的暴露分配方式的效应。例如，考虑只针对社区中连接最多的前10\%成员的干预 $\bm g_2^*$。与完全随机地向10\%的社区成员分配暴露相比，预计这种有针对性的干预，对整个社区维持健身房会员资格的平均概率会有更大的影响。假想干预 $\bm g_3^*$ 为朋友数少于10的人额外分配一位进行身体活动的朋友。值得注意的是，该干预旨在直接干预部分社区成员的社交网络。最后，估计在同一社区中同时实施动态干预 $\bm g_2^*$ 和基于网络的干预 $\bm g_3^*$ 的联合效应。这里报告每种干预下的结局期望；定义为这些干预之间对比的因果效应，也可以用我们的方法估计，其结果见补充材料。

模拟分别采用样本量 $n=500$、$n=1000$ 和 $n=10{,}000$，每个数据集均按上述数据生成分布抽样。在同一网络的条件下，抽取1000个这样的数据集，重复进行估计；对每种样本量，网络只抽取一次。每个单位 $i$ 的基线协变量按独立同分布抽样。我们用 TMLE 估计上述干预 $\bm g_1^*$ 至 $\bm g_4^*$ 下结局期望的样本平均，并计算相应估计值真实的有限样本偏差与方差。

我们还通过模拟，比较 \texttt{tmlenet} R 包提供的三种不同渐近方差估计量。具体而言，比较它们相应的渐近置信区间的覆盖率。第一种，是把观测当作独立同分布数据，使用有效影响曲线方差的独立同分布代入估计量（“IID Var”）。第二种，是对相关观测作正确调整的代入方差估计量（“dependent IC Var”）。最后一种，是前面介绍的参数自助方差估计量（“bootstrap Var”）。图21.3按干预类型、置信区间类型以及观测样本量分层，给出95\%置信区间的平均长度和覆盖率。

我们通过95\%置信区间的覆盖率，评估上述三种方差估计方法的性能。结果表明，在这类数据集中忽略相依性而开展推断，通常会产生偏于不保守的方差估计，并使置信区间覆盖不足；即使样本量非常大，覆盖率也可能低至50\%，见图21.3右侧面板中的“IID Var”。当样本量足够大时，基于相依方差估计的置信区间可获得接近名义95\%的覆盖率，见同图右侧面板中的“dependent IC Var”；但在小样本中，由于尚未达到渐近正态性以及接近违反正值性，其表现可能不佳。值得注意的是，基于参数自助方差估计的置信区间，在较小样本量下提供了最稳健的覆盖率，并且在大样本下，几乎所有模拟情境均达到名义95\%覆盖率，见同图右侧面板中的“bootstrap Var”。

\begin{landscape}
\begin{figure}[p]
 \centering
 % 原书真实矢量数据层；文字在数据区域之外替换。不得用模拟数据重建。
\begin{tikzpicture}[x=0.59mm,y=-0.59mm]
\node[anchor=north west,inner sep=0pt] at (0,0) {\includegraphics[width=199.4200mm,height=142.7800mm]{figures/ch21_simulation_data.pdf}};
\node[font=\fontsize{9.5}{10.5}\selectfont,inner sep=0pt] at (96.4849,9.0394) {区间类型};
\node[font=\fontsize{9.5}{10.5}\selectfont,inner sep=0pt,anchor=west] at (123.2980,8.9302) {相依 IC 方差};
\node[font=\fontsize{9.5}{10.5}\selectfont,inner sep=0pt,anchor=west] at (182.0741,8.9302) {自助法方差};
\node[font=\fontsize{9.5}{10.5}\selectfont,inner sep=0pt,anchor=west] at (230.8036,8.9302) {独立同分布方差};
\node[font=\fontsize{9.5}{10.5}\selectfont,inner sep=0pt] at (142.7137,235.9497) {均值估计与95\%置信区间长度};
\node[font=\fontsize{9.5}{10.5}\selectfont,inner sep=0pt] at (272.8118,236.3217) {覆盖率};
\node[anchor=east,font=\fontsize{9.5}{10.5}\selectfont,inner sep=0pt] at (70.8000,32.5464) {$g_1^*$（随机35\%）};
\node[anchor=east,font=\fontsize{9.5}{10.5}\selectfont,inner sep=0pt] at (70.8000,47.8160) {$g_2^*$（动态干预）};
\node[anchor=east,font=\fontsize{9.5}{10.5}\selectfont,inner sep=0pt] at (70.8000,63.0797) {$g_3^*$（网络干预）};
\node[anchor=east,font=\fontsize{9.5}{10.5}\selectfont,inner sep=0pt] at (70.8000,78.3496) {$g_2^*+g_3^*$};
\node[anchor=east,font=\fontsize{9.5}{10.5}\selectfont,inner sep=0pt] at (70.8000,100.4004) {$g_1^*$（随机35\%）};
\node[anchor=east,font=\fontsize{9.5}{10.5}\selectfont,inner sep=0pt] at (70.8000,115.6635) {$g_2^*$（动态干预）};
\node[anchor=east,font=\fontsize{9.5}{10.5}\selectfont,inner sep=0pt] at (70.8000,130.9339) {$g_3^*$（网络干预）};
\node[anchor=east,font=\fontsize{9.5}{10.5}\selectfont,inner sep=0pt] at (70.8000,146.2036) {$g_2^*+g_3^*$};
\node[anchor=east,font=\fontsize{9.5}{10.5}\selectfont,inner sep=0pt] at (70.8000,168.2479) {$g_1^*$（随机35\%）};
\node[anchor=east,font=\fontsize{9.5}{10.5}\selectfont,inner sep=0pt] at (70.8000,183.5178) {$g_2^*$（动态干预）};
\node[anchor=east,font=\fontsize{9.5}{10.5}\selectfont,inner sep=0pt] at (70.8000,198.7879) {$g_3^*$（网络干预）};
\node[anchor=east,font=\fontsize{9.5}{10.5}\selectfont,inner sep=0pt] at (70.8000,214.0510) {$g_2^*+g_3^*$};
\end{tikzpicture}

 \caption{优先连接网络中，按样本量、干预和置信区间类型分组的平均95\%置信区间长度（左侧面板）及覆盖率（右侧面板）。仅给出平均结局期望的结果。}
 \label{fig:ch21:simulation}
\end{figure}
\end{landscape}

本研究一个出人意料的发现，是参数自助法在小样本推断中表现出明显的稳健性。具体而言，预料之中的是，这些模拟所使用的高度相连网络类型，可能导致由相依数据影响曲线方差估计推导出的渐近置信区间覆盖率偏于不保守，例如由于接近违反正值性和尚未收敛至正态分布。然而，基于参数自助方差估计的类似置信区间，即使样本量低至500，仍给出了接近名义95\%的覆盖率，这一点令人意外。未来研究将探索参数自助法发挥作用所需的假设，以及其对违反这些假设的敏感性。

此外，模拟说明，与相同水平、无针对性的暴露覆盖相比，通过有针对性地向连接密切且进行身体活动的个体分配暴露，可以利用网络结构，提高维持健身房会员资格的平均概率。我们还展示了，直接对观测到的网络结构施加干预，例如干预 $\bm g_3^*$，其效应是可以估计的；这类干预还能与经济激励结合，例如干预 $\bm g_2^*+\bm g_3^*$。这种联合干预可能对整个社区产生更大的平均期望效应，尤其是在干预受到资源约束时。

\section{注释与延伸阅读}

关于网络及网络数据中因果推断的文献正在迅速发展。不过，目前用于网络中因果效应估计与推断的统计方法仍然有限，这一领域的文献也只是最近才开始发展（van der Laan，2014a；VanderWeele 等，2014b；Ogburn 和 VanderWeele，2014；VanderWeele 和 An，2013；Tchetgen Tchetgen 和 VanderWeele，2012）。近期提出的大多数方法可分为两类：一类依赖于单位之间的暴露随机化假设（Rosenbaum，2007；Aronow 和 Samii，2013；Bowers 等，2013；Walker 和 Muchnik，2014；Aral 和 Walker，2011、2014；Toulis 和 Kao，2013；Liu 和 Hudgens，2014；Choi，2014；Basse 和 Airoldi，2015），另一类则对结局进行参数建模，将其表示为单位网络的特定函数。

在网络情境中使用过的一些参数方法，包括广义线性模型（generalized linear model，GLM）与广义估计方程（generalized estimating equation，GEE）（Christakis 和 Fowler，2013、2007）；这些方法具有重要局限（Lyons，2010；VanderWeele，2011、2013；VanderWeele 等，2012c；Ogburn 和 VanderWeele，2014）。首先，GLM 及类似的建模技术，要求对底层数据生成过程作出较强的简化建模假设。因此，网络数据情境中 GEE 与 GLM 的模型误设，是一个主要问题。更为关键的也许是，使用 GLM 和其他类似统计技术开展有效统计推断，通常要求观测单位相互独立，而由于网络数据本身的性质，这一假设不太可能成立。此前也已有研究指出，把这类标准统计程序用于相依数据，会导致无效且通常偏于不保守的统计推断（Lyons，2010；Ogburn 和 VanderWeele，2014）。

此外，近年来也开始出现一些很有希望的网络数据估计方法。例如，Aronow 和 Samii（2013）在随机化研究情境下提出 Horvitz--Thompson 估计量，定义所谓“网络暴露模型”，并推导有限样本方差估计量。但这些方法在观察性研究情境中的用途有限。其他用于识别和估计网络中处理效应的方法，包括随机行动者导向模型（stochastic actor-oriented model；Steglich 等，2010），以及能够容纳网络不确定性的线性贝叶斯建模方法（Toulis 和 Kao，2013）。近期提出的另一种方法，把靶向极大似然估计（targeted maximum likelihood estimation）的半参数框架应用于观察性网络数据情境（van der Laan，2014a），在允许更大且更现实的一类数据生成模型的同时，得到有效的渐近推断。
